% 论文创新点

本文的主要创新点如下：

\begin{enumerate}
  \item \textbf{编码器感知的强化学习前处理框架}。针对NVENC等硬件编码器"不可导"的固有限制，本文提出将其作为强化学习中的"黑盒环境"纳入优化回路。通过Actor-Critic架构，系统以NVENC编码后的实际码率和感知质量（LPIPS）作为奖励信号，使前处理网络RAPNet能够在不依赖编码器梯度的前提下，自适应地学习针对硬件编码器特性的预补偿策略，有效缓解了低码率场景下的质量崩塌问题。

  \item \textbf{基于Random Crop的高频保持训练策略}。区别于传统的下采样训练方式（如2$\times$下采样到540p），本文采用从1080p原始帧中随机裁剪540$\times$960区域的训练策略。该策略在保持与下采样方案相同计算开销的前提下，完整保留了原始视频的高频纹理和边缘细节，避免了下采样导致的信息丢失，使模型能够学习到更精细的增强策略。

  \item \textbf{频域感知的复合损失函数设计}。除策略梯度损失外，本文设计了一套基于空间频率分解的辅助损失体系：（1）低频结构保持损失（通过$5 \times 5$均值滤波约束低频内容一致性）；（2）非边缘区域高频抑制损失（利用Sobel边缘检测生成软掩模，仅在平坦区域惩罚高频修改，避免边缘模糊）；（3）边缘聚焦损失（鼓励模型将修改集中在边缘区域，提升主观锐度）。该设计使模型在提升编码效率的同时保持了视觉自然度。
\end{enumerate}